% TEMPLATE-MILC-v3.tex - plantilla maestra UNICA de las guias MILC v3.
%
% La usan las tres familias de guias, cada una con su driver:
%   · Grados 6-11  -> scripts/build-guias-g11.py
%   · Bebras       -> scripts/build-guias-bebras.py
%   · Semillero    -> scripts/build-guias-semillero.py
%
% Antes habia tres copias casi identicas (~400 lineas iguales, 6 distintas):
% cada mejora habia que aplicarla tres veces, y en la practica no se hacia
% (el motor de recursos vivia solo en dos de ellas). Lo que varia entre
% familias esta parametrizado en 6 placeholders que cada driver llena
% (se nombran aqui SIN su sintaxis real: el builder reemplaza por texto plano,
% tambien dentro de los comentarios, y un valor multilinea rompe el preambulo):
%
%   HEADER_IZQ        encabezado izquierdo de pagina
%   HEADER_DER        encabezado derecho de pagina
%   PORTADA_ETIQUETA  rotulo sobre el numero grande (GRADO/MOMENTO/MODULO)
%   PORTADA_SERIE     linea de serie de la portada
%   PORTADA_ID        identificador de la guia en la portada
%   PORTADA_CIRCULO   el circulito de grado (°), o vacio si no aplica
%
% El resto de claves las documenta content/guias/_SCHEMA.md. El script valida
% que no queden placeholders sin reemplazar antes de escribir el archivo final.

\documentclass[11pt,letterpaper]{article}
\usepackage[margin=2.0cm]{geometry}
\usepackage[table]{xcolor}
\usepackage{fontspec}
\usepackage[spanish]{babel}
\usepackage{tikz}
% Librerías TikZ para los diagramas de las guías (bloque `recursos:`):
%  · babel  → babel en español vuelve activos caracteres como «>», y TikZ los usa
%             en sus opciones (>=stealth, ->). Sin esta librería, cualquier
%             diagrama con flechas revienta con «\language@active@arg> has an
%             extra }». Debe ir ANTES de que se dibuje nada.
%  · positioning        → sintaxis «below left=2pt and 3pt».
%  · decorations.pathreplacing → llaves ({brace}) para señalar rangos.
\usetikzlibrary{babel,positioning,decorations.pathreplacing}
\usepackage{graphicx}
% Circuitos: el trabajo de electrónica se hace hoy con micro:bit y sus sensores
% integrados (MakeCode, sin cableado), así que no se necesitan esquemáticos.
% Si algún día entran montajes con protoboard, basta descomentar esta línea:
% los diagramas .tex/.tikz declarados en `recursos:` ya se incrustan solos.
% \usepackage{circuitikz}
\usepackage[most]{tcolorbox}
\usepackage{tabularx}
\usepackage{array}
\usepackage{fancyhdr}
\usepackage{titlesec}
\usepackage[document]{ragged2e}  % bandera a la derecha en todo el cuerpo
\usepackage{enumitem}
\usepackage{microtype}

% Helvetica Neue no trae la flecha U+2192, asi que cada «→» escrito literalmente
% en el YAML se compilaba a NADA, en silencio (462 flechas en 61 guias: rutas del
% tipo «paleta → tipografia → lockup» salian rotas). Mapeamos el caracter a su
% version matematica, que si tiene glifo. Asi el autor puede seguir escribiendo
% «→» comodamente en el YAML.
\usepackage{newunicodechar}
\newunicodechar{→}{\ensuremath{\rightarrow}}
% Mismo problema con los simbolos de marca y vinetas: el «✓» de «una palomita ✓»
% salia como cuadrito vacio en el PDF de todas las guias que lo usaban.
\usepackage{pifont}
\newunicodechar{✓}{\ding{51}}
\newunicodechar{✗}{\ding{55}}
\newunicodechar{✔}{\ding{52}}
\newunicodechar{•}{\textbullet}
\newunicodechar{≥}{\ensuremath{\geq}}
\newunicodechar{≤}{\ensuremath{\leq}}

\setmainfont{Helvetica Neue}
\setsansfont{Helvetica Neue}
\setlength{\parindent}{0pt}
\setlength{\parskip}{6pt}
% Sin particion de palabras: en 6.º-7.º una palabra cortada por guion al final
% de linea («Comuni-cacion») frena la lectura mucho mas que una linea corta.
\hyphenpenalty=10000
\exhyphenpenalty=10000
\pretolerance=2000
\tolerance=3000
\setlength{\emergencystretch}{3em}
\linespread{1.10}
\renewcommand{\arraystretch}{1.25}
\setlist{leftmargin=5mm,itemsep=1mm,topsep=1mm}
\newcolumntype{Y}{>{\RaggedRight\arraybackslash}X}

\definecolor{milcMagenta}{HTML}{E6007E}
\definecolor{milcVino}{HTML}{5A0038}
\definecolor{milcTurquesa}{HTML}{0093A5}
\definecolor{milcVerde}{HTML}{58A923}
\definecolor{milcMostaza}{HTML}{E5B400}
\definecolor{milcOcre}{HTML}{B86600}
\definecolor{milcNegro}{HTML}{241816}
\definecolor{milcGris}{HTML}{F7F7F4}
\definecolor{milcLinea}{HTML}{E8E2DD}
\definecolor{milcGradePrimary}{HTML}{0D9488}
\definecolor{milcGradeDark}{HTML}{115E59}
\definecolor{milcGradeSoft}{HTML}{CCFBF1}
\definecolor{milcGradeAccent}{HTML}{CFE7FF}
\definecolor{milcGradeText}{HTML}{FFFFFF}
\color{milcNegro}

\pagestyle{fancy}
\fancyhf{}
\lhead{\small Territorio interior · Educación socioemocional}
\rhead{\small Grado 9° · Momento 3 · MILC v3}
\cfoot{\small\thepage}
\renewcommand{\headrulewidth}{0pt}

\newtcolorbox{titlebox}[1]{
  enhanced,colback=#1,colframe=#1,arc=7mm,boxrule=0pt,
  left=7mm,right=7mm,top=3mm,bottom=3mm,
  before skip=8pt,after skip=8pt,
  fontupper=\sffamily\bfseries\Large\color{white}
}

\newtcolorbox{softbox}[3]{
  enhanced,colback=#2,colframe=#1,borderline west={4pt}{0pt}{#1},
  arc=2mm,boxrule=.4pt,left=4mm,right=4mm,top=3mm,bottom=3mm,
  before skip=6pt,after skip=8pt,title={#3},coltitle=white,
  colbacktitle=#1,fonttitle=\sffamily\bfseries,fontupper=\RaggedRight,
  attach boxed title to top left={xshift=3mm,yshift=-2mm},
  boxed title style={arc=2mm, boxrule=0pt}
}

\newtcolorbox{drawbox}[2]{
  enhanced,colback=white,colframe=#1,arc=4mm,boxrule=1pt,
  left=5mm,right=5mm,top=4mm,bottom=4mm,before skip=8pt,after skip=8pt,
  title={#2},coltitle=white,colbacktitle=#1,fonttitle=\sffamily\bfseries,
  fontupper=\RaggedRight,
  attach boxed title to top left={xshift=4mm,yshift=-2mm},
  boxed title style={arc=3mm, boxrule=0pt}
}

\newtcolorbox{infoband}[1]{
  enhanced,colback=#1!8,colframe=#1!50,arc=2mm,boxrule=.5pt,
  left=4mm,right=4mm,top=2mm,bottom=2mm,before skip=4pt,after skip=4pt,
  fontupper=\small\RaggedRight
}

% Puente narrativo entre secciones (contrato editorial Conéctate).
% Da continuidad: "Acabas de X. Ahora vas a Y porque Z."
% Visualmente: banda izquierda en color del grado, texto en itálica pequeña.
\newtcolorbox{puentebox}{
  enhanced,colback=milcGradeSoft,colframe=milcGradeDark,
  borderline west={3pt}{0pt}{milcGradeDark},
  arc=0mm,boxrule=0pt,left=5mm,right=4mm,top=2.5mm,bottom=2.5mm,
  before skip=4pt,after skip=8pt,
  fontupper=\itshape\small\color{milcNegro}\RaggedRight
}
% La flecha va en modo matematico a proposito: Helvetica Neue NO tiene el glifo
% U+2192, asi que \textrightarrow se compilaba a NADA («Missing character: There
% is no → in font Helvetica Neue») y el puente salia sin flecha. En modo
% matematico la flecha viene de la fuente matematica y si se dibuja.
\newcommand{\puente}[1]{%
  \begin{puentebox}%
    \textcolor{milcGradeDark}{$\rightarrow$}\hspace{2mm}#1%
  \end{puentebox}%
}

\newcommand{\linead}[1]{\noindent\color{milcLinea}\rule{#1}{0.5pt}\color{milcNegro}}
\newcommand{\respuesta}[1]{\par\vspace{2mm}\foreach \n in {1,...,#1}{\noindent\color{milcLinea}\rule{\linewidth}{0.5pt}\par\vspace{5mm}}\color{milcNegro}}
\newcommand{\checkbox}{\textcolor{milcGradeDark}{\fbox{\phantom{x}}}\hspace{5pt}}

% ─── Listas aireadas para las guías socioemocionales ────────────────────
% Componer las verificaciones como «\checkbox texto\\» deja el texto que salta
% de línea debajo del cuadrito y apelmaza el bloque. Estos entornos alinean la
% continuación y airean los ítems. Los usa Territorio Interior; cualquier
% familia puede hacerlo.
\newlist{checklista}{itemize}{1}
\setlist[checklista]{
  label=\textcolor{milcGradeDark}{\fbox{\phantom{x}}},
  leftmargin=9mm, labelsep=3.2mm, itemsep=2.4mm,
  topsep=2.5mm, parsep=0pt, partopsep=0pt, before=\RaggedRight
}
\newlist{pasoslista}{enumerate}{1}
\setlist[pasoslista]{
  label=\textcolor{milcGradeDark}{\sffamily\bfseries\arabic*.},
  leftmargin=8mm, labelsep=3mm, itemsep=2.8mm,
  topsep=1mm, parsep=0pt, partopsep=0pt, before=\RaggedRight
}

\newcommand{\phasecard}[4]{
\begin{tcolorbox}[enhanced,colback=#3,colframe=#2,arc=3mm,boxrule=.5pt,height=3.05cm,valign=center,halign=center,left=1.5mm,right=1.5mm,top=2mm,bottom=2mm]
{\sffamily\bfseries\fontsize{22}{24}\selectfont\textcolor{#2}{#1}}\par
{\sffamily\bfseries\small #4}
\end{tcolorbox}
}

% ── Piezas del rediseno 2026 (legibilidad 6.º-11.º) ───────────────────
% Banda de estacion: reemplaza el titlebox macizo. Titulo a la izquierda,
% dato practico (tiempo/modalidad) a la derecha, en una sola linea.
\newcommand{\milcestacion}[2]{%
\begin{tcolorbox}[enhanced,colback=#1,colframe=#1,arc=2mm,boxrule=0pt,
  left=5mm,right=5mm,top=2.6mm,bottom=2.6mm,before skip=11pt,after skip=7pt]
{\sffamily\bfseries\fontsize{15}{18}\selectfont\color{white}#2}
\end{tcolorbox}}

% Tarjeta de paso numerada: sustituye las tablas «Pilar / Aplicacion», que
% eran formato de consulta docente, no de estudiante. El numero va en un
% circulo sobre el borde para que la secuencia se lea de un vistazo.
\newcommand{\milcpaso}[3]{%
\begin{tcolorbox}[enhanced,colback=white,colframe=milcLinea,arc=1.5mm,boxrule=.9pt,
  left=14mm,right=4mm,top=3mm,bottom=3mm,before skip=4pt,after skip=4pt,
  fontupper=\RaggedRight,
  overlay={\node[anchor=north west,circle,fill=#2,text=white,inner sep=0pt,
    minimum size=7.4mm,font=\sffamily\bfseries\small]
    at ([xshift=3.6mm,yshift=-3.2mm]frame.north west) {#1};}]
#3
\end{tcolorbox}}

% Casilla marcable de tamano util con lapiz (la anterior era \fbox{\phantom{x}},
% de alto variable segun el texto que la seguia).
\newcommand{\milcmarca}{\framebox[3.8mm]{\rule{0pt}{3.4mm}}}

% Fila marcable de lista suelta (checks de la estacion 1).
\newcommand{\milccheckrow}[1]{%
\par\vspace{1.8mm}\noindent
\makebox[6.4mm][l]{\raisebox{-.55mm}{\textcolor{milcNegro}{\milcmarca}}}%
\parbox[t]{\dimexpr\linewidth-6.4mm\relax}{\RaggedRight #1}\par}

% Celda marcable dentro de la rubrica (tabla de autoevaluacion). Misma
% sangria francesa, pero pensada para vivir dentro de una columna X.
\newcommand{\milcopcion}[1]{%
\makebox[5.2mm][l]{\raisebox{-.55mm}{\textcolor{milcNegro}{\milcmarca}}}%
\parbox[t]{\dimexpr\linewidth-5.2mm\relax}{\RaggedRight #1}}

% ── Recursos visuales de la guía ──────────────────────────────────────
% Declarados en el bloque `recursos:` del YAML y emitidos por el builder.
% \guiaFigura   → imagen rasterizada o PDF (foto, carta celeste, montaje).
% \guiaDiagrama → fuente TikZ/circuitikz (.tex) incrustada como vector:
%                 es la vía para circuitos y esquemáticos editables.
% #1 = ruta relativa al .tex · #2 = pie (puede ir vacío)
\newcommand{\guiaFigura}[2]{%
\begin{center}
\includegraphics[width=0.88\linewidth,height=0.38\textheight,keepaspectratio]{#1}\\[1.5mm]
{\footnotesize\itshape\textcolor{milcNegro!75}{#2}}
\end{center}
}

\newcommand{\guiaDiagrama}[2]{%
\begin{center}
\input{#1}\\[1.5mm]
{\footnotesize\itshape\textcolor{milcNegro!75}{#2}}
\end{center}
}

\begin{document}
\thispagestyle{empty}

% ── Portada ───────────────────────────────────────────────────────────
% Antes: un rectangulo de color a pagina completa. Se veia bien en pantalla,
% pero en la fotocopiadora del colegio se comia el toner y salia gris sucio.
% Ahora la banda ocupa el tercio superior y el resto va en blanco.
\begin{tikzpicture}[remember picture,overlay]
  \path[fill=milcGradePrimary]
    (current page.north west) rectangle ([yshift=-10.6cm]current page.north east);

  \node[anchor=north west,text=milcGradeText] at ([xshift=2.0cm,yshift=-1.55cm]current page.north west)
    {\sffamily\bfseries\fontsize{12}{14}\selectfont MOMENTO};
  \node[anchor=north west,text=milcGradeText,opacity=.85] at ([xshift=2.0cm,yshift=-2.35cm]current page.north west)
    {\sffamily\bfseries\fontsize{8.5}{10}\selectfont TERRITORIO INTERIOR · SOCIOEMOCIONAL};
  \node[anchor=north east,text=milcGradeText,opacity=.85,align=right] at ([xshift=-2.0cm,yshift=-1.55cm]current page.north east)
    {\sffamily\bfseries\fontsize{9}{12}\selectfont I.E. Sor María Juliana\\Cartago, Valle del Cauca};

  \node[anchor=west,text=milcGradeText] (gradonum) at ([xshift=1.85cm,yshift=-7.1cm]current page.north west)
    {\sffamily\bfseries\fontsize{112}{112}\selectfont 3};
  % El circulito de grado (°) solo aplica a las guias de grado («11°»); en
  % Bebras y Semillero el numero grande es un momento/modulo, no un grado.
  % Se posiciona relativo al nodo `gradonum`, no en coordenadas absolutas.
  

  \node[anchor=west,text=milcGradeText,text width=11.4cm,align=left] at ([xshift=1.5cm,yshift=0.5cm]gradonum.east)
    {
     {\sffamily\bfseries\fontsize{9}{11}\selectfont\MakeUppercase{Territorio interior · Socioemocional}}\\[3mm]
     {\sffamily\bfseries\fontsize{21}{25}\selectfont\setlength{\baselineskip}{27pt}Vínculos y\\[4pt]consentimiento\\[4pt]quién decide de verdad}\\[3.5mm]
     {\sffamily\bfseries\fontsize{15}{18}\selectfont Grado 9° · Momento 3}
    };
\end{tikzpicture}

\vspace*{9.4cm}
\vfill

{\sffamily\bfseries\fontsize{8}{10}\selectfont\textcolor{milcGradeDark}{LO QUE TE LLEVAS HOY}}

\begin{tcolorbox}[enhanced,colback=white,colframe=milcNegro,arc=2mm,boxrule=1.6pt,
  left=5mm,right=5mm,top=4mm,bottom=4mm,before skip=3pt,after skip=12pt,
  fontupper=\RaggedRight\fontsize{11}{15.5}\selectfont]
Tu lista de señales tempranas: las cinco libertades revisadas en una relación real —propia o cercana—, la distinción escrita entre cuidado y control con ejemplos de cada uno, y la ruta con nombre y lugar de a quién acudirías.
\end{tcolorbox}

\begin{tcolorbox}[enhanced,colback=milcGris,colframe=milcGris,arc=2mm,boxrule=0pt,
  left=5mm,right=5mm,top=3.5mm,bottom=3.5mm,after skip=12pt]
{\sffamily\bfseries\fontsize{8}{10}\selectfont\textcolor{milcNegro!55}{ESTA GUÍA ES DE}}\\[3mm]
\begin{tabularx}{\linewidth}{X p{3.2cm} p{3.2cm}}
\linead{\linewidth} & \linead{3.0cm} & \linead{3.0cm}\\[-1mm]
{\fontsize{8.5}{10}\selectfont\textcolor{milcNegro!55}{Nombre completo}} &
{\fontsize{8.5}{10}\selectfont\textcolor{milcNegro!55}{Grupo}} &
{\fontsize{8.5}{10}\selectfont\textcolor{milcNegro!55}{Fecha}}\\
\end{tabularx}
\end{tcolorbox}

\vfill

\noindent\textcolor{milcLinea}{\rule{\linewidth}{1pt}}\\[2mm]
{\sffamily\bfseries\fontsize{9.5}{12}\selectfont Autor: Dr. Álvaro Cárdenas Orozco}\\[1mm]
{\sffamily\fontsize{8.5}{11}\selectfont\textcolor{milcNegro!55}{Metodología MILC · Apertura ancestral · Vínculos y consentimiento · Triángulo Dussel-Estoicismo-Floridi}}

\newpage

\section*{Apertura: Vínculos y consentimiento: quién decide, y las señales que llegan temprano}

\begin{softbox}{milcGradeDark}{milcGradeSoft}{Saber ancestral}
Hubo un tiempo, y no tan lejano, en que \textbf{la decisión sobre con quién unir la vida no era de los dos que se unían}. El \textbf{23 de marzo de 1776}, Carlos III sancionó una \textbf{Pragmática Sanción} con un título que lo dice todo: «\emph{para evitar el abuso de contraer matrimonios desiguales}». \textbf{Obligaba a los hijos de familia menores de veinticinco años a \emph{pedir y obtener} el consentimiento del padre} para casarse. \textbf{Se extendió a América} por real cédula dada en El Pardo el \textbf{7 de abril de 1778}, y estuvo tan viva que se modificó y repitió una y otra vez durante décadas, hasta ser prácticamente resancionada el \textbf{1.º de junio de 1803}. \emph{Fíjate en la palabra clave del título: \textbf{desiguales}.} \textbf{La ley no se preocupaba por el amor: se preocupaba por que la gente no se uniera cruzando las líneas sociales} ---y en la América colonial esas líneas eran también de color, como vas a ver en el momento 5---. \textbf{Guarda entonces el punto:} \emph{durante siglos, el consentimiento sobre un vínculo \textbf{le perteneció a alguien que no estaba en el vínculo}}. \emph{Hoy ninguna ley dice eso. Y sin embargo, en muchas relaciones, la decisión sigue sin ser de los dos.}
\end{softbox}

\begin{softbox}{milcTurquesa}{milcGris}{Saber contemporáneo}
¿Cómo se le quita a alguien la decisión sobre su propia vida sin que nadie lo note? \textbf{El investigador Evan Stark lo estudió y le puso un nombre: \emph{control coercitivo}.} Lo describe combinando dos cosas: la \textbf{coerción} ---usar fuerza o amenazas para provocar o impedir una respuesta--- y el \textbf{control} ---formas de privación y de mando que consiguen obediencia \textbf{de manera indirecta}---. \textbf{Cuando las dos van juntas se produce lo que él llama una \emph{condición de no libertad}, que se vive como estar atrapado} (Stark, 2007). \textbf{Y aquí está su aporte decisivo, el que cambia cómo se mira todo esto:} \emph{lo que ocurre en muchas relaciones dañinas \textbf{no es principalmente violento ni ocurre en episodios}} ---es \textbf{un patrón sostenido de conductas de control}---. Por eso Stark propone llamarlo un \textbf{delito contra la libertad}. \textbf{Y describe cómo se hace, en cosas pequeñísimas:} \emph{controlar con quién se habla, revisar el teléfono, decidir la ropa, exigir respuesta inmediata, llevar cuenta de las horas, ir apartando a la persona de sus amigos}. \emph{Ninguna de esas cosas, sola, parece grave. Esa es exactamente la razón por la que funciona.}
\end{softbox}

\begin{softbox}{milcMostaza}{milcGris}{Pregunta-puente}
\textit{En 1776 el permiso sobre un vínculo era del padre, \textbf{y estaba escrito}. \textbf{¿Qué decisiones tuyas ---con quién hablas, qué te pones, a qué hora contestas--- están hoy en manos de otra persona} \ldots \textbf{sin que nada esté escrito}?}
\end{softbox}

\begin{softbox}{milcVerde}{milcGris}{Saber y saber hacer}
Al terminar podrás: (1) \textbf{entender} que el consentimiento sobre los vínculos fue durante siglos de otros, y qué significa que hoy sea tuyo; (2) \textbf{explicar} el consentimiento como un proceso ---libre, informado, específico y \textbf{revocable}--- y no como un permiso que se da una vez; (3) \textbf{distinguir} cuidado de control usando el criterio del patrón y el de la libertad; (4) \textbf{revisar} cinco libertades concretas en una relación y \textbf{saber a quién acudir}, con nombre y lugar.
\end{softbox}



\small
\begin{tabularx}{\textwidth}{>{\bfseries}p{3.0cm}Y}
\rowcolor{milcGradePrimary!12}Autor & Dr. Álvaro Cárdenas Orozco\\
DBA & Comprende los fenómenos que afectan a su territorio, regula sus emociones ante ellos y acompaña a otros con empatía (transversal 6°--11°, articulado con la política «Escuela, territorio de vida» del MEN, Resolución 006519 de 2025, y la Ley 1620 de 2013)\\
Referentes & Primera ayuda psicológica (OMS, 2011/2012) · Directrices IASC (2007) · USGS y Servicio Geológico Colombiano · Pueblo nasa y la minga (CRIC) · Dussel · Estoicismo · Floridi\\
Producto final & Tu lista de señales tempranas: las cinco libertades revisadas en una relación real —propia o cercana—, la distinción escrita entre cuidado y control con ejemplos de cada uno, y la ruta con nombre y lugar de a quién acudirías.\\
\end{tabularx}
\normalsize

\puente{En el momento 2 viste lo que se espera de cada uno. \textbf{Hoy vas a lo que pasa cuando esas expectativas entran en una relación de pareja} ---donde el mandato se disfraza de amor---.}

\milcestacion{milcGradeDark}{Ruta MILC: así vamos a aprender}
\begin{tabularx}{\textwidth}{>{\centering\arraybackslash}X>{\centering\arraybackslash}X>{\centering\arraybackslash}X>{\centering\arraybackslash}X}
\phasecard{1}{milcVerde}{milcGris}{Escucha\\[-1mm]\small Observo, escucho y pregunto.} &
\phasecard{2}{milcTurquesa}{milcGris}{Sistematización\\[-1mm]\small Distingo cuidado de control.} &
\phasecard{3}{milcMagenta}{milcGris}{Praxis\\[-1mm]\small Construyo y aplico.} &
\phasecard{4}{milcVino}{milcGris}{Evaluación\\[-1mm]\small Reflexión liberadora.}
\end{tabularx}


\puente{Primero, sin hablar de nadie: qué se dice en tu curso que es «normal» en una relación. \emph{Vas a encontrar cosas que no lo son.}}

\milcestacion{milcVerde}{Estación 1 de 3 · Escucha}

\begin{softbox}{milcVerde}{milcGris}{Escena para escuchar}
\textbf{Actividad 1 · IDENTIFICA --- Lo que se dice que es normal.} \textbf{Sin nombres, nunca. Ni de tu curso ni de tu casa.} \textbf{Primera parte.} Escribe \textbf{las frases que circulan} sobre las relaciones de pareja en tu edad: «si no te cela es porque no le importas», «es que es muy celoso pero es que me quiere mucho», «pásame tu clave si no tienes nada que esconder», «no me gusta que hables con él». \emph{Escríbelas tal cual se dicen.} \textbf{Segunda parte.} Al frente de cada una, responde una sola pregunta: \emph{¿esto le \textbf{quita} algo a alguien?} \textbf{Tiempo, amigos, privacidad, decisiones, tranquilidad.} \textbf{Tercera parte.} Piensa en \textbf{una relación que conozcas} ---de lejos, sin nombres--- y escribe \textbf{qué podía hacer esa persona antes de la relación que ya no hace}. \textbf{Y la última, para ti solo:} ¿le has pedido a alguien, o te han pedido, \textbf{la clave del celular}? \emph{Escribe qué se dijo y cómo se sintió.}
\end{softbox}

{\sffamily\bfseries\fontsize{8}{10}\selectfont\textcolor{milcNegro!55}{MARCA CUANDO LO HAYAS HECHO}}
\milccheckrow{Escribiste las frases tal como circulan, sin corregirlas.}
\milccheckrow{Respondiste para cada una si \textbf{le quita algo} a alguien.}
\milccheckrow{Escribiste qué dejó de hacer una persona que conoces desde que está en esa relación.}

\begin{infoband}{milcVerde}


\textbf{En tu cuaderno (Actividad 1 · IDENTIFICA).} Título: «Actividad 1. Lo que se dice que es normal». \textbf{Formato:} las frases con la pregunta respondida + lo que alguien dejó de hacer + tu propia respuesta sobre la clave. \textbf{Extensión:} media página. \textbf{Sabes que terminaste cuando} encontraste \textbf{al menos una frase que suena a amor y funciona quitando algo}. Esta página es tuya: \textbf{nadie más tiene que leerla si tú no quieres}.
\end{infoband}

\puente{Ya lo tienes. Ahora, qué es consentir de verdad, por qué el control es un patrón y no un episodio, y cuáles son las señales tempranas.}

\milcestacion{milcTurquesa}{Fase 2 · Sistematización: entender el consentimiento}

\begin{softbox}{milcTurquesa}{milcGris}{Cuatro claves sobre consentir}
Cuatro claves para distinguir lo que cuida de lo que controla, cuando las dos cosas se dicen con las mismas palabras.
\end{softbox}

\milcpaso{1}{milcTurquesa}{\textbf{Durante siglos el permiso fue de otro.} La Pragmática de 1776 obligaba a los menores de veinticinco a \textbf{obtener el consentimiento del padre}, y se aplicó en América desde 1778. \emph{Su objetivo declarado eran los matrimonios «desiguales»:} \textbf{no cuidaba a nadie ---cuidaba un orden social---.} \textbf{Y esto sirve para desarmar un argumento que vas a oír toda la vida:} \emph{«siempre ha sido así»}. \textbf{No: cambió, y cambió hace poco.} \emph{Que hoy la decisión sea de los dos que están en el vínculo es una conquista reciente ---y por eso hay que saber en qué consiste, porque lo que se conquistó también se puede ir perdiendo de a poquitos.}}
\milcpaso{2}{milcTurquesa}{\textbf{Consentir es un proceso, no una firma.} Un consentimiento que valga algo tiene cuatro condiciones: es \textbf{libre} ---sin presión, sin chantaje, sin insistencia hasta que el otro se rinda---; es \textbf{informado} ---la persona sabe a qué está diciendo que sí---; es \textbf{específico} ---decir que sí a una cosa no es decir que sí a otra---; y es \textbf{revocable} ---se puede cambiar de opinión en cualquier momento, incluso después de haber dicho que sí---. \emph{Esa última condición es la que más se ignora y la que más importa.} \textbf{Y hay una consecuencia directa:} \emph{la insistencia no produce consentimiento}. \textbf{Un «bueno, ya» dicho para que el otro pare no es un sí ---es una rendición---.}}
\milcpaso{3}{milcTurquesa}{\textbf{El control es un patrón, no un episodio.} Este es el aporte de Stark y cambia por completo dónde hay que mirar: \emph{lo que hace daño en muchas relaciones \textbf{no es principalmente violento ni ocurre en episodios sueltos}} ---es un \textbf{patrón sostenido} de coerción y control que produce una \textbf{condición de no libertad} (Stark, 2007). \textbf{Por eso preguntar «¿te ha pegado?» es la pregunta equivocada, o al menos la tardía.} \emph{La pregunta correcta es otra:} \textbf{«¿qué podías hacer antes que ya no haces?»}. \emph{Y por eso también nadie recuerda un momento en que la cosa empezó:} \textbf{no empezó ---se fue estrechando---}, y cada paso, mirado solo, parecía razonable.}
\milcpaso{4}{milcTurquesa}{\textbf{Las cinco libertades: el examen concreto.} Stark describe cómo la persona controlada va perdiendo cosas específicas: \emph{salir sin pedir permiso, ver a sus amigos, decir lo que piensa, ponerse lo que quiere, disponer de su propio tiempo}. \textbf{Conviértelas en preguntas y tienes un examen que cualquiera puede aplicar:} \emph{¿puede salir sin avisar y sin que haya problema? ¿sigue viendo a sus amigos de antes? ¿dice lo que piensa delante de la otra persona? ¿se pone lo que quiere? ¿su tiempo es suyo?} \textbf{Y la pregunta que las resume:} \emph{¿esta relación le hizo la vida \textbf{más grande} o \textbf{más pequeña}?} \textbf{El cuidado agranda la vida de alguien. El control la encoge, y siempre en nombre del amor.}}

\begin{drawbox}{milcTurquesa}{Cuidado o control: la misma frase, dos cosas distintas}
\begin{pasoslista}
\item \textbf{«Avísame que llegaste»} — \emph{cuidado} si es tranquilidad; \textbf{control} si es rendir cuentas y hay problema si no.
\item \textbf{«No me gusta que hables con él»} — \emph{decir lo que a uno le incomoda} es válido; \textbf{prohibirlo} no lo es.
\item \textbf{«Muéstrame el celular»} — \emph{no hay versión de cuidado.} \emph{La confianza no se verifica: si hay que revisar, ya no hay.}
\item \textbf{«Es que me pongo mal si no contestas»} — \emph{contar lo que uno siente} es válido; \textbf{usarlo para obligar} es chantaje.
\item \textbf{Prueba final:} \emph{¿le hizo la vida más grande o más pequeña?}
\end{pasoslista}
\end{drawbox}

\begin{softbox}{milcMostaza}{milcGris}{Errores comunes y su corrección}
\textbf{«Si no te cela es porque no le importas»:} confunde una emoción con una conducta. \textbf{Buscar solo golpes:} el patrón empieza mucho antes y casi nunca ahí. \textbf{«Pero si tiene sus cosas buenas»:} todas las relaciones las tienen; eso no anula el patrón. \textbf{Pedir la clave «si no tienes nada que esconder»:} la privacidad no es sospecha. \textbf{Dar consejos de salida a un amigo:} rara vez funcionan; acompañar y avisar a un adulto, sí. \textbf{«Ya le dije que cambie»:} el patrón no se corrige con avisos. \textbf{Pensar que es cosa de mujeres:} le pasa a cualquiera, y a los hombres les cuesta más contarlo por el mandato del momento 2.
\end{softbox}

\begin{infoband}{milcTurquesa}


\textbf{En tu cuaderno (Actividad 2 · EXPLICA).} Título: «Actividad 2. Cuidado o control». \textbf{Formato:} toma \textbf{cuatro} frases de tu Actividad 1 y escribe, para cada una, \textbf{la versión que cuida} y \textbf{la versión que controla}. \textbf{Extensión:} media página. \textbf{Sabes que terminaste cuando} puedes explicar por qué \textbf{«¿te ha pegado?» es la pregunta tardía} y cuál es la temprana.
\end{infoband}

\puente{Ya sabes distinguir. Es hora de aplicarlo ---a una relación real, propia o cercana--- y de tener la ruta lista antes de necesitarla.}

\milcestacion{milcMagenta}{Fase 3 · Praxis: revisar las señales}

\begin{softbox}{milcMagenta}{milcGris}{Cómo se revisan las señales tempranas}
Esto se practica sobre casos, no sobre personas del salón. Y termina donde tiene que terminar: en saber a quién acudir antes de necesitarlo.
\end{softbox}

\milcpaso{1}{milcMagenta}{\textbf{Cuidado o control (10 min, en parejas).} Trabajen con \textbf{casos inventados} ---nunca con alguien del colegio, ni siquiera sin nombre--- y clasifiquen diez conductas. \emph{Discutan las que no coincidan:} \textbf{el desacuerdo casi siempre está en las conductas que \emph{una sola vez} parecen razonables}, y ahí es donde se aprende, porque el criterio no es la conducta suelta sino \textbf{si se repite y si quita algo}.}
\milcpaso{2}{milcMagenta}{\textbf{Las cinco libertades (8 min).} Escribe las cinco preguntas y aplícalas a \textbf{una relación} ---la tuya, si tienes; si no, una que conozcas de lejos, sin nombres---. \emph{Y hazte la pregunta que las resume:} \textbf{¿la vida de esa persona se hizo más grande o más pequeña?}}
\milcpaso{3}{milcMagenta}{\textbf{La línea de tiempo (8 min).} Toma un caso ---inventado o conocido de lejos--- y \textbf{ordena las conductas por orden de aparición}: qué pasó primero, qué después. \emph{Van a ver el patrón que describe Stark:} \textbf{empieza por lo que parece cariño ---escribir todo el día, querer estar siempre--- y avanza hacia lo que quita}. \emph{Ver el orden es lo que permite reconocerlo temprano la próxima vez.}}
\milcpaso{4}{milcMagenta}{\textbf{La ruta, otra vez (7 min).} \textbf{Escribe a quién acudirías}, con \textbf{nombre, cargo y lugar} ---orientación escolar, un docente de confianza, un adulto de tu casa---. \emph{Es la misma ruta del momento 5 del grado 8, y aquí aplica igual:} \textbf{si hay amenazas, contenido sexual o riesgo físico, eso no se maneja entre amigos.} \textbf{Y agrega dos frases:} una para \textbf{pedir ayuda para ti} y otra para \textbf{acompañar a alguien} ---«te creo, no es normal, no estás solo, vamos juntos a hablar con \_\_\_\_»---.}

\begin{drawbox}{milcMagenta}{Antes de cerrar tu lista de señales}
\begin{checklista}
\item Trabajaste con casos inventados o lejanos: \textbf{ningún nombre de nadie del colegio}.
\item Escribiste, para cuatro frases, la versión que cuida y la que controla.
\item Aplicaste \textbf{las cinco libertades} y respondiste si la vida se hizo más grande o más pequeña.
\item Tu línea de tiempo muestra que \textbf{el patrón empieza por lo que parece cariño}.
\item Tu ruta tiene \textbf{nombre, cargo y lugar}.
\item Tienes escritas las dos frases: para pedir ayuda y para acompañar.
\end{checklista}
\end{drawbox}

\begin{softbox}{milcMagenta}{milcGris}{Plantilla de trabajo}
\textbf{Las cinco libertades.} \emph{«¿Sale sin pedir permiso? · ¿Ve a sus amigos? · ¿Dice lo que piensa? · ¿Se pone lo que quiere? · ¿Su tiempo es suyo?»} \\[1mm]
\textbf{La pregunta que resume.} \emph{«¿Esta relación le hizo la vida más grande o más pequeña?»} \\[1mm]
\textbf{Para acompañar.} \emph{«Te creo. Eso no es normal. No estás solo. ¿Vamos juntos a hablar con \_\_\_\_?»}
\end{softbox}

\begin{infoband}{milcMagenta}


\textbf{En tu cuaderno (Actividad 3 · CREA).} Título: «Actividad 3. Señales tempranas». \textbf{Formato:} las cuatro frases en sus dos versiones + las cinco libertades aplicadas + la línea de tiempo del patrón + la ruta con nombre y lugar y las dos frases. \textbf{Extensión:} una página. \textbf{Sabes que terminaste cuando} tu ruta tiene \textbf{un nombre propio}, no «un profesor».
\end{infoband}

\puente{El producto incluye una ruta con nombre propio. \emph{Una ruta que hay que buscar en el momento no sirve; ya lo aprendiste en el grado 8.}}

\milcestacion{milcMostaza}{Producto de la sesión}
\textbf{Lista de señales tempranas: cuidado frente a control, las cinco libertades revisadas y la ruta con nombre y lugar}.
\begin{softbox}{milcMostaza}{milcGris}{Criterios de calidad}
(1) El trabajo se hace sobre casos inventados o lejanos, sin identificar a nadie del colegio. (2) Cada frase aparece en sus dos versiones —la que cuida y la que controla— con el criterio explicado. (3) Las cinco libertades están aplicadas y se responde si la vida de la persona se amplió o se redujo. (4) La línea de tiempo muestra el orden del patrón, empezando por lo que parece cariño. (5) La ruta identifica persona, cargo y lugar, y reconoce qué situaciones no se manejan entre pares.
\end{softbox}

\puente{Antes de cerrar, revisa lo más importante de esta guía: si tus ejemplos de control son todos golpes, te faltaron justamente las señales tempranas.}

\milcestacion{milcVino}{Evaluación liberadora · cinco dimensiones}

\begin{softbox}{milcVino}{milcGris}{Sentido de la evaluación}
Evaluar no es solo revisar una nota. Es reconocer qué aprendí, qué cambió en mí, qué emociones manejé, cómo me relaciono con la ciudad y la comunidad, y qué le dejaré a quien viene detrás.
\end{softbox}

\textbf{Mi reflexión liberadora en cinco dimensiones.} Completa cada frase iniciada:

\small
\begin{tabularx}{\textwidth}{>{\bfseries\RaggedRight\arraybackslash}p{3.4cm}Y>{\RaggedRight\arraybackslash}p{4.6cm}}
\rowcolor{milcVino}\color{white}Dimensión & \color{white}\bfseries Mi respuesta (frame starter) & \color{white}\bfseries Algo que descubrí\\
1. Personal & Lo que más cambió en mí fue \linead{6.5cm} & \linead{4.2cm}\\
2. Emocional & La emoción más fuerte fue \linead{2.5cm} y la regulé \linead{3cm} & \linead{4.2cm}\\
3. Ciudadana & En democracia esto importa porque \linead{6cm} & \linead{4.2cm}\\
4. Local (Cartago) & En mi barrio sería útil para \linead{6.5cm} & \linead{4.2cm}\\
5. Intergeneracional & A mi abuela/abuelo le diría \linead{6.5cm} & \linead{4.2cm}\\
\end{tabularx}
\normalsize

\begin{infoband}{milcVino}
\textbf{Para la próxima sustentación.} Vuelve a esta tabla en seis meses. Verás que las respuestas cambian — no porque lo aprendido cambie, sino porque tú cambias. La evaluación liberadora es una conversación contigo a través del tiempo.
\end{infoband}


\puente{Cerramos con tres voces: la comunidad que se funda escuchando, lo que está en tu poder ---incluido el rechazo--- y lo que cambia cuando todo queda registrado.}

\milcestacion{milcGradeDark}{Triángulo de pensamiento · cierre filosófico}

\begin{softbox}{milcVino}{milcGris}{Dussel · lente del nosotros}
\textit{«Aquel que es capaz de escuchar silenciosamente al Otro como otro es el que puede constituir una comunidad y no una mera sociedad totalizada.»}

{\footnotesize\itshape\color{milcNegro!70}--- Enrique Dussel · Filosofía de la liberación (1977)}

Fíjate en tres palabras: \textbf{al Otro \emph{como otro}}. \emph{Es decir: como alguien distinto, con vida propia, que no es una extensión de uno.} \textbf{Y eso es exactamente lo que el control no soporta:} \emph{quien controla no quiere a la otra persona como otra ---la quiere como parte suya---}, y por eso le molestan sus amigos, su tiempo, sus opiniones, su ropa. \textbf{Todo lo que en ella no viene de él le sobra.} \emph{Y aquí está el criterio más limpio de toda la guía:} \textbf{una relación buena \emph{agranda} las dos vidas} ---cada uno conoce más gente, hace más cosas, se atreve a más---. \emph{Si dos vidas se hicieron más pequeñas, eso no era una totalidad amorosa: era una totalidad, y ya.}

\textbf{En las relaciones que veo alrededor, ¿las dos personas hacen más cosas que antes, o menos?}
\end{softbox}
\linead{\linewidth}\\[1mm]
\linead{\linewidth}

\begin{softbox}{milcMostaza}{milcGris}{Estoicismo · Epicteto · Enquiridión, 1 (c. 125 d.C.; trad. desde E. Carter) · cuidado interior}
\textit{«Algunas cosas están en nuestro poder y otras no. Están en nuestro poder la opinión, el impulso, el deseo, el rechazo; en una palabra, todo aquello que es obra nuestra.»}

{\footnotesize\itshape\color{milcNegro!70}--- Epicteto · Enquiridión, 1 (c. 125 d.C.; trad. desde E. Carter)}

Otra vez la palabra \textbf{rechazo}, y aquí es literal: \emph{decir que no ---a algo, a un plan, a un contacto, a una foto--- \textbf{es obra tuya} y no deja de serlo por estar en una relación}. \textbf{Y hay una segunda lectura, más incómoda, dirigida al que siente celos:} \emph{el celo es una emoción y llega solo ---eso no está en tu poder---}. \textbf{Lo que sí es obra tuya es \emph{qué haces} con él:} \emph{revisar un teléfono, exigir una respuesta, prohibir una amistad}. \textbf{Sentir celos no le hace daño a nadie. Actuarlos sí.} \emph{Y esa distinción es la misma que trabajaste en el momento 2 del grado 7, ahora con mucho más en juego.}

\textbf{Lo último que hice por celos, ¿fue una emoción que sentí o una conducta que elegí?}
\end{softbox}
\linead{\linewidth}\\[1mm]
\linead{\linewidth}

\begin{softbox}{milcTurquesa}{milcGris}{Floridi · lente de la infoesfera}
\textit{«Las tecnologías digitales están re-ontologizando la naturaleza misma de la infoesfera.»}

{\footnotesize\itshape\color{milcNegro!70}--- Luciano Floridi · La cuarta revolución (2014; trad.)}

El control coercitivo que describe Stark necesitaba antes \textbf{presencia}: para saber dónde estaba alguien había que ir. \emph{Hoy no.} \textbf{La ubicación compartida, el doble visto, la hora de última conexión y el «¿por qué no contestaste?» convierten el control en algo que se puede ejercer \emph{a toda hora y desde lejos}, y encima con apariencia de normalidad.} \emph{Eso es re-ontologizar:} \textbf{no es que el control sea más frecuente ---es que cambió de naturaleza---.} \emph{Y hay una consecuencia práctica que conviene decir sin rodeos:} \textbf{compartir ubicación permanente y claves con una pareja no es una prueba de amor} ---es la infraestructura exacta que el control necesita para funcionar---.

\textbf{¿Alguien puede saber en todo momento dónde estoy? ¿Se lo ofrecí yo o me lo pidieron?}
\end{softbox}
\linead{\linewidth}\\[1mm]
\linead{\linewidth}

\begin{infoband}{milcNegro}
\textbf{Nota para el docente.} Las tres son citas textuales y llevan su referencia debajo. Puedes remitir a la obra en clase.
\end{infoband}

\milcestacion{milcVino}{Mi autoevaluación · marca una casilla por fila}

{\small
\setlength{\extrarowheight}{2.2pt}
\begin{tabularx}{\textwidth}{>{\RaggedRight\arraybackslash\bfseries}p{3.0cm}YYY}
\rowcolor{milcVino}\color{white}\bfseries Criterio & \color{white}\bfseries Lo logré & \color{white}\bfseries En proceso & \color{white}\bfseries Necesito apoyo\\[.6mm]
Comprensión del tema & \milcopcion{Explico las ideas con mis palabras.} & \milcopcion{Reconozco las ideas, falta precisión.} & \milcopcion{Necesito repasar lo básico.}\\[.8mm]
\rowcolor{milcGris}Calidad del análisis & \milcopcion{Distingo cuidado de control con ejemplos, y sé a quién acudir con nombre y lugar.} & \milcopcion{Reconozco las señales cuando son graves, pero no las tempranas.} & \milcopcion{Sigo pensando que los celos y el control son formas de querer mucho.}\\[.8mm]
Aplicación y producto & \milcopcion{Mi producto cumple los criterios.} & \milcopcion{Está armado, con ajustes pendientes.} & \milcopcion{Aún está en construcción.}\\[.8mm]
\rowcolor{milcGris}Reflexión 5 dimensiones & \milcopcion{Respondí cada una con honestidad.} & \milcopcion{Respondí algunas con detalle.} & \milcopcion{Mis respuestas son muy generales.}\\[.8mm]
Triángulo de pensamiento & \milcopcion{Conecté las 3 voces con mi trabajo.} & \milcopcion{Conecté al menos una.} & \milcopcion{No las relaciono con mi proceso.}\\[.8mm]
\end{tabularx}}

\vspace{3mm}
\textbf{Hoy entendí que:}\\[1mm]
\linead{\linewidth}\\[1mm]
\linead{\linewidth}

\textbf{Mi compromiso para la próxima sesión es:}\\[1mm]
\linead{\linewidth}\\[1mm]
\linead{\linewidth}

\begin{softbox}{milcMostaza}{milcGris}{Cita de despedida · Marco Aurelio}
\textit{``El obstáculo en el camino se vuelve el camino. Lo que estorba al camino se vuelve camino.''}\\[1mm]
Cada error de hoy, bien observado, se convierte en pieza del camino que pisarás mañana.
\end{softbox}

\small
\begin{tabularx}{\textwidth}{>{\bfseries}p{3.6cm}Y}
\rowcolor{milcGradeDark!12}Nombre completo: & \linead{10cm}\\
Fecha: & \linead{4cm} \quad Curso/grupo: \linead{4cm}\\
Mi firma: & \linead{9cm}\\
\end{tabularx}
\normalsize

\begin{infoband}{milcGradeDark}
\textbf{Para la próxima sesión.} Dos cosas. \textbf{Primero, deja tu ruta escrita donde puedas verla} ---en el cuaderno, en una nota del celular--- y, si te alcanza el valor, \textbf{preséntate con la persona de orientación}. \emph{No hace falta contar nada: basta con haber hablado una vez, para que pedir ayuda después no sea empezar de cero.} \textbf{Y si al hacer este ejercicio reconociste algo en tu propia relación, o en la de alguien cercano: habla con un adulto.} \emph{No es exagerar. Reconocerlo temprano es exactamente para lo que sirve esta guía.} \textbf{Segundo, pregunta en tu casa cómo se decidía antes:} averigua con alguien mayor \textbf{cómo se conocieron sus padres o sus abuelos}, \textbf{si alguien tuvo que dar permiso}, si alguna vez una familia se opuso a un noviazgo y \textbf{por qué razón se oponía}. \textbf{Trae:} la ruta escrita y lo que te contaron. \emph{Vas a encontrar la Pragmática de 1776 viva en las historias de tu familia: casi siempre la oposición no era por la persona, sino por lo que representaba ---de dónde venía, de qué color era, cuánto tenía---.}
\end{infoband}

\end{document}
